\section{Introduction}

\label{sec:photon-introduction}

In standard metastable consensus, a validator processes one transaction at a time per
conflict set: it samples, waits for responses, updates confidence, and repeats for
$\beta$ rounds. This creates a throughput ceiling of $1/(\beta \cdot \Delta)$ finalizations
per second per conflict set, where $\Delta$ is the round time. For $\beta = 11$ and
$\Delta = 50$ms, this is merely 1.8 finalizations per second per conflict set.

While the Nova protocol amortizes this cost across DAG subgraphs, independent conflict
sets still bottleneck at the per-conflict finalization rate. Photon addresses this by
pipelining: multiple transactions enter the sampling pipeline simultaneously, sharing
each round's query-response cycle across all in-flight transactions.

The key challenge is ensuring that pipelining does not compromise safety. When multiple
transactions are sampled in the same round, their confidence counters may interact
(e.g., via shared Byzantine responses). We prove that this interaction is benign: the
confidence dynamics for independent conflict sets are statistically independent, and
for transactions within the same conflict set, pipelining reduces to sequential processing
of the conflict resolution.

\paragraph{Contributions.}
\begin{enumerate}[nosep]
\item The Photon pipelined consensus protocol with confidence matrix formalization (\S\ref{sec:photon-protocol}).
\item A proof that pipelined safety degrades at most linearly in pipeline depth (\S\ref{sec:photon-safety}).
\item Pipeline stall analysis and resolution bounds (\S\ref{sec:photon-stalls}).
\item Throughput analysis showing sustained finalization at the transaction arrival rate (\S\ref{sec:photon-throughput}).
\item Empirical evaluation achieving 10,000 TPS (\S\ref{sec:photon-evaluation}).
\end{enumerate}

%% -----------------------------------------------------------------------
%%  2. SYSTEM MODEL
%% -----------------------------------------------------------------------
